\subsection{Порядок выполнения}
\textbf{Первый проход:} проследите путь одного дома от строки таблицы через подготовку признаков до прогноза цены. Выполните B0 и объясните MAE в денежных единицах. К сравнению десяти конфигураций переходите после проверки этого запуска.

\subsubsection*{Шаг 1. Изучите исходные таблицы}
Откройте обе таблицы и описание данных. Проверьте число строк и столбцов,
наличие \file{SalePrice} только в обучающем файле и уникальность
\file{Id} внутри каждого файла. Назовите два числовых и два
категориальных признака; объясните их смысл.

\begin{lstlisting}
from pathlib import Path
import pandas as pd

DATA_DIR = Path("data")
train_data = pd.read_csv(DATA_DIR / "train.csv")
test_data = pd.read_csv(DATA_DIR / "test.csv")
print(train_data.shape, test_data.shape)
print(train_data.head())
print(test_data.head())
print(train_data.isna().sum().sort_values(ascending=False).head(10))
\end{lstlisting}

\begin{keybox}{Контрольная точка 1}
В \file{train.csv} цена известна, в \file{test.csv} --- отсутствует.
\file{Id} нужен для связи прогноза с домом и не включается в признаки.
Пропущенную целевую цену нельзя заполнять медианой: такая строка не даёт
достоверного обучающего ответа.
\end{keybox}

\subsubsection*{Шаг 2. Выполните разделение и предобработку}
Сохраните фиксированное удаление пяти столбцов из исходного задания.
Отделите \file{SalePrice} и \file{Id}, затем разделите сырые строки
в отношении 80/20 с \file{random\_state=40}. Только теперь обучайте
преобразования. Готовая реализация находится в \file{lab2\_utils.py}.
Прочитайте функцию \file{prepare\_frames} и найдите в ней место,
где вычисляются обучающие статистики.

\begin{lstlisting}
import sys
sys.path.insert(0, "code")
from lab2_utils import load_course_data, constant_baselines

data = load_course_data(DATA_DIR, seed=40)
print(data["Atr"].shape, data["Ava"].shape, data["Atest"].shape)
print(data["removed_columns"])
print(data["feature_names"][:10])
\end{lstlisting}

Числовой конвейер объединяет заполнение медианой и стандартизацию.
Категориальный --- заполнение отдельной категорией и one-hot-кодирование.
\file{ColumnTransformer} применяет эти конвейеры к соответствующим
столбцам~\cite{columntransformer}. Признаки с целочисленными кодами
проверьте по словарю: тип хранения сам по себе не доказывает,
что признак количественный.

\begin{keybox}{Контрольная точка 2}
Три матрицы имеют одинаковое число столбцов и не содержат NaN или
бесконечностей. Пересечение \file{train\_ids} и
\file{validation\_ids} пусто. Объекты валидации не участвовали
в \file{fit} преобразований. Число входов сети равно
\file{data["Atr"].shape[1]}.
\end{keybox}

\subsubsection*{Шаг 3. Рассчитайте простые базовые прогнозы}
Для сравнения используйте константные прогнозы: среднюю и медианную цену
\textit{обучающей} части. Для каждого вычислите MAE и RMSE на валидации.

\begin{lstlisting}
baseline_table = pd.DataFrame(constant_baselines(data)).T
print(baseline_table)
\end{lstlisting}

Такая проверка показывает, даёт ли сложная модель пользу по сравнению
с решением, не использующим признаки. Нейросеть не обязана выигрывать:
если это не произошло, разберите причины и честно покажите результат.

\subsubsection*{Шаг 4. Постройте исходную нейросеть}
Создайте и проанализируйте последовательную модель. Исходная настройка:
один скрытый слой из 64 нейронов с ReLU, один линейный выход, Adam,
скорость обучения \(10^{-3}\), MSE. Входная размерность определяется
результатом предобработки.

\begin{lstlisting}
from tensorflow import keras

keras.utils.set_random_seed(40)
d = data["Atr"].shape[1]
model = keras.Sequential([
    keras.Input(shape=(d,)),
    keras.layers.Dense(64, activation="relu"),
    keras.layers.Dense(1, activation="linear"),
])
model.compile(
    optimizer=keras.optimizers.Adam(learning_rate=1e-3),
    loss="mse",
    metrics=[keras.metrics.MeanAbsoluteError(name="mae")],
)
model.summary()
\end{lstlisting}

Этот небольшой пример показывает работу API. Основное самостоятельное
задание --- изменить архитектуру, провести предусмотренные исследования
и объяснить результаты. Для своей реализации рассчитайте число параметров
вручную и проверьте его по сводке модели.

\subsubsection*{Шаг 5. Обучите модель и сохраните историю}
Передавайте в модель \file{Atr} и стандартизованную цель \file{ztr},
а валидационные данные задавайте явно. Не используйте переменную
\file{test\_data} в аргументе \file{validation\_data}.

\begin{lstlisting}
history = model.fit(
    data["Atr"], data["ztr"],
    validation_data=(data["Ava"], data["zva"]),
    epochs=100, batch_size=32, shuffle=True, verbose=0,
)
\end{lstlisting}

Сопровождающая функция \file{run\_experiment} объединяет создание
новой модели, обучение, обратное преобразование прогноза и расчёт метрик.
Используйте либо собственную последовательность шагов, либо эту функцию;
не обучайте две модели подряд по ошибке, принимая их за один опыт.

\begin{lstlisting}
from lab2_utils import run_experiment, plot_diagnostics

result = run_experiment(
    data, hidden=(64,), epochs=100, batch_size=32,
    optimizer="adam", learning_rate=1e-3,
    loss="mse", seed=40, early_stop=False,
)
print(result["metrics"])
fig = plot_diagnostics(data, result)
\end{lstlisting}

\subsubsection*{Шаг 6. Оцените качество и исследуйте ошибки}
Постройте три графика:
\begin{enumerate}
\item MAE обучения и валидации в зависимости от эпохи.
\item Предсказанная цена против фактической на валидации; добавьте прямую
\(\widehat y=y\), обозначающую идеальный прогноз.
\item Остаток \(e=y-\widehat y\) в зависимости от предсказанной цены;
добавьте горизонтальную линию \(e=0\).
\end{enumerate}
Во всех подписях укажите используемые единицы. По последним двум графикам
обсудите крупные промахи, систематическое завышение или занижение цены,
а также изменение разброса ошибки.

Выделите пять объектов с наибольшей абсолютной ошибкой и рассмотрите
их исходные характеристики. Не удаляйте их из валидации ради улучшения
метрики. Если предполагаете ошибку данных, сформулируйте проверяемую
гипотезу и укажите, каких сведений не хватает.

\begin{keybox}{Контрольная точка 3}
MAE и RMSE вычислены после обратного преобразования цели. Для каждой
конфигурации использованы те же валидационные объекты. У каждой кривой
указано, к какой модели и части данных она относится.
\end{keybox}

\subsubsection*{Шаг 7. Получите и сохраните прогноз}
После выбора конфигурации используйте её для всех строк
\file{test.csv}. В базовом варианте сохраняется модель, обученная на
80\% размеченных строк. Её прогнозы сопровождаются оценкой, полученной
на оставшихся 20\%.

\begin{lstlisting}
import numpy as np
from lab2_utils import save_experiment

preds = np.asarray(result["predictions"]).reshape(-1)
output = pd.DataFrame({
    "Id": data["prediction_ids"],
    "SalePrice": preds,
})
assert len(output) == len(test_data)
assert output["Id"].tolist() == test_data["Id"].tolist()
assert np.isfinite(output["SalePrice"]).all()
print("Negative prices:", int((output["SalePrice"] < 0).sum()))
save_experiment(data, result, "runs/baseline_01")
\end{lstlisting}

Не обрезайте отрицательные прогнозы незаметно. Сначала проверьте обучение
и обратное преобразование. Если выбирается правило постобработки,
обоснуйте его по валидации и одинаково примените к обоим наборам
прогнозов. Файл \file{predictions\_raw.csv} сохраняет исходные значения.

Допускается отдельное финальное переобучение на всём \file{train.csv}
после фиксации настроек. В этом случае нужно заново обучить
предобработку и сеть. Прежняя валидационная ошибка относится к прежней
модели: вычислять её на строках, включённых в финальное обучение,
как независимую оценку нельзя.

\Needspace{20\baselineskip}
\subsection{Выполнение на PyTorch}
Вы можете выполнить работу на PyTorch, сохранив те же разбиение,
предобработку, архитектуру, метрики и исследуемые факторы.
В обычном интерфейсе PyTorch обучение не выполняется вызовами
\file{model.compile()} и \file{model.fit()}: цикл записывается явно.
Минимальная последовательность действий внутри цикла~\cite{pytorch}:

\begin{lstlisting}
model.train()
for xb, yb in train_loader:
    optimizer.zero_grad()
    prediction = model(xb)
    loss = criterion(prediction, yb)
    loss.backward()
    optimizer.step()

model.eval()
with torch.no_grad():
    val_prediction = model(X_val_tensor)
\end{lstlisting}

В этом фрагменте загрузчик, модель, оптимизатор и тензоры требуется
создать самостоятельно. Проверьте, что прогноз и цель имеют одинаковую
форму \((b,1)\). Сочетание \((b,1)\) и \((b,)\) может привести
к нежелательному расширению размерностей при вычислении потерь.
В качестве функций потерь используйте \file{MSELoss} и \file{L1Loss}.
Результаты TensorFlow и PyTorch не обязаны совпадать при одинаковом seed:
сравнивайте соблюдение протокола и полученные метрики, а не побитовое
совпадение весов.
